\documentclass[10pt,a4paper]{article}
\usepackage[utf8]{inputenc}
\usepackage[T1]{fontenc}
\usepackage[ngerman]{babel}
\usepackage{geometry}
\usepackage{xcolor}
\usepackage{tabularx}
\usepackage{array}
\usepackage{booktabs}
\usepackage{enumitem}
\usepackage{helvet}
\usepackage{tikz}
\usepackage{colortbl}
\usepackage{amssymb}
\usepackage{fancyhdr}
\usepackage{titlesec}
\usepackage{tocloft}
\usepackage{hyperref}

\usetikzlibrary{shapes.geometric, arrows.meta, positioning, calc}

\renewcommand{\familydefault}{\sfdefault}
\geometry{margin=2cm, top=2.5cm, bottom=2.5cm}
\setlength{\parindent}{0pt}
\setlength{\parskip}{6pt}
\setlength{\headheight}{14pt}

% Colors
\definecolor{spanisch}{HTML}{c41e3a}
\definecolor{lightred}{HTML}{fdf2f4}
\definecolor{darkgray}{HTML}{333333}
\definecolor{lightgray}{HTML}{f5f5f5}
\definecolor{spaet}{HTML}{8e44ad}
\definecolor{lightspaet}{HTML}{f5eef8}
\definecolor{success}{HTML}{27ae60}
\definecolor{warning}{HTML}{f39c12}

% Headers
\pagestyle{fancy}
\fancyhf{}
\fancyhead[L]{\small\textcolor{spaet}{\textbf{Kompendium Spanisch Spätbeginner -- Mecklenburg-Vorpommern}}}
\fancyhead[R]{\small\textcolor{darkgray}{v2.0 | Januar 2026}}
\fancyfoot[C]{\thepage}

% Section formatting
\titleformat{\section}{\Large\bfseries\color{spaet}}{}{0pt}{}[\vspace{-4pt}\textcolor{spaet}{\rule{\textwidth}{1pt}}]
\titleformat{\subsection}{\large\bfseries\color{darkgray}}{}{0pt}{}
\titleformat{\subsubsection}{\normalsize\bfseries\color{spanisch}}{}{0pt}{}
\titlespacing{\section}{0pt}{16pt}{8pt}
\titlespacing{\subsection}{0pt}{12pt}{4pt}
\titlespacing{\subsubsection}{0pt}{8pt}{2pt}

% Hyperref setup
\hypersetup{
    colorlinks=true,
    linkcolor=spaet,
    urlcolor=spanisch
}

% Custom commands
\newcommand{\wichtig}[1]{%
\begin{center}
\fcolorbox{spanisch}{lightred}{\parbox{0.92\textwidth}{\small #1}}
\end{center}
}

\newcommand{\hinweis}[1]{%
\begin{center}
\fcolorbox{spaet}{lightspaet}{\parbox{0.92\textwidth}{\small #1}}
\end{center}
}

\newcommand{\quelle}[1]{\textcolor{darkgray}{\scriptsize [Quelle: #1]}}

\begin{document}

%% TITLE PAGE
\begin{titlepage}
\begin{center}
\vspace*{2cm}

{\Huge\bfseries\textcolor{spaet}{KOMPENDIUM}}\\[0.5cm]
{\LARGE\textcolor{darkgray}{Spanisch als neu begonnene Fremdsprache}}\\[0.3cm]
{\LARGE\textcolor{darkgray}{(Spätbeginner)}}\\[1cm]

{\Large Mecklenburg-Vorpommern}\\[0.5cm]
{\large Klassenstufen 10--12 | Gesamtschule \& Gymnasium}\\[2cm]

\begin{tikzpicture}[scale=1.2]
\node[draw, fill=lightspaet, minimum width=1.8cm, minimum height=1cm, rounded corners=2pt] (a1) at (0,0) {\textbf{A1}};
\node[draw, fill=lightspaet, minimum width=1.8cm, minimum height=1cm, rounded corners=2pt] (a2) at (3,0) {\textbf{A2}};
\node[draw, fill=lightspaet, minimum width=1.8cm, minimum height=1cm, rounded corners=2pt] (a2p) at (6,0) {\textbf{A2+}};
\node[draw, fill=spaet, text=white, minimum width=1.8cm, minimum height=1cm, rounded corners=2pt] (b1) at (9,0) {\textbf{B1(+)}};

\draw[-{Stealth[length=3mm]}, thick, color=spaet] (a1) -- (a2);
\draw[-{Stealth[length=3mm]}, thick, color=spaet] (a2) -- (a2p);
\draw[-{Stealth[length=3mm]}, thick, color=spaet] (a2p) -- (b1);

\node[below] at (0,-0.8) {\footnotesize Start};
\node[below] at (3,-0.8) {\footnotesize Klasse 10};
\node[below] at (6,-0.8) {\footnotesize Klasse 11};
\node[below] at (9,-0.8) {\footnotesize Abitur};
\end{tikzpicture}

\vspace{2cm}

\textbf{Für die Lehrkräfte-Anerkennung und Unterrichtsplanung}\\[0.5cm]
{\small Enthält: Rechtliche Grundlagen, Unterrichtsplanung nach Klassenstufe,\\
Prüfungsformate, Lehrwerk-Analyse, praktische Hinweise}

\vfill

{\small Version 2.0 | Januar 2026}\\
{\footnotesize Lernpfade | Spanisch Spätbeginner}

\end{center}
\end{titlepage}

%% TABLE OF CONTENTS
\tableofcontents
\newpage

%% ============================================================================
%% TEIL 1: GRUNDLAGEN
%% ============================================================================

\section*{TEIL 1: GRUNDLAGEN}
\addcontentsline{toc}{section}{TEIL 1: GRUNDLAGEN}

\subsection{Rechtlicher Rahmen und Definition}

\subsubsection{Was bedeutet ``neu begonnene Fremdsprache''?}

Der Begriff \textbf{``Spätbeginner''} ist eine informelle Bezeichnung. Die offiziellen Dokumente des Landes Mecklenburg-Vorpommern verwenden den Terminus:

\wichtig{
\textbf{``In der Einführungsphase neu begonnene Fremdsprache''}\\[4pt]
Definition: Spanisch ab Klassenstufe 10 (Einführungsphase) ohne Vorkenntnisse.\\
Es existiert \textbf{kein separater Rahmenplan} für diese Kursform.
}

\subsubsection{Die maßgeblichen Rechtsquellen}

\begin{center}
\small
\begin{tabular}{|p{4.5cm}|p{4cm}|p{5cm}|}
\hline
\rowcolor{lightgray}
\textbf{Dokument} & \textbf{Herausgeber} & \textbf{Relevanz} \\
\hline
\textbf{Rahmenplan Spanisch Qualifikationsphase} (2019) & Ministerium für Bildung M-V & Verbindliche Themenfelder, Kompetenzen \\
\hline
\textbf{Oberstufen- und Abiturprüfungsverordnung} (2019) & Land M-V & Stundentafel, Prüfungsformate, Klausurdauer \\
\hline
Handreichung Sprechprüfung (2020) & Institut für Qualitätsentwicklung & Format, Bewertungsbögen \\
\hline
Handreichung Paarprüfung Abitur (2025) & Institut für Qualitätsentwicklung & Mündliche Abiturprüfung \\
\hline
Kommentierte Literaturempfehlungen & Institut für Qualitätsentwicklung & Lektürevorschläge (nicht verbindlich) \\
\hline
\end{tabular}
\end{center}

\subsubsection{Neu begonnene Fremdsprache = eigene Kategorie}

\wichtig{
\textbf{Zentrale Erkenntnis für die Unterrichtsplanung:}\\[4pt]
Die neu begonnene Fremdsprache ist \textbf{weder} Grundkurs (3 Wochenstunden) \textbf{noch} Leistungskurs (5 Wochenstunden), sondern eine \textbf{eigenständige Kategorie} mit \textbf{4 Wochenstunden}.\\[4pt]
Laut Oberstufen- und Abiturprüfungsverordnung: \textit{``Eine neu beginnende Fremdsprache wird \textbf{vierstündig} unterrichtet.''}
}

\begin{center}
\small
\begin{tabular}{|l|c|c|c|}
\hline
\rowcolor{lightgray}
\textbf{Kursart} & \textbf{Wochenstunden} & \textbf{Leistungskurs möglich?} & \textbf{Zielniveau Abitur} \\
\hline
Grundkurs (fortgeführte Fremdsprache) & 3 & Nein & B2 \\
\hline
Leistungskurs (fortgeführte Fremdsprache) & 5 & Ja & B2+ \\
\hline
\rowcolor{lightspaet}
\textbf{Neu begonnene Fremdsprache} & \textbf{4} & \textbf{Nein} & \textbf{B1(+)} \\
\hline
\end{tabular}
\end{center}

\quelle{Oberstufen- und Abiturprüfungsverordnung M-V (2019)}

\subsection{Gemeinsamer Europäischer Referenzrahmen: Die Progression}

\subsubsection{Zielniveaus für Spätbeginner}

\begin{center}
\begin{tikzpicture}[
    node distance=0.8cm,
    box/.style={draw, fill=lightspaet, minimum width=2.5cm, minimum height=1.2cm, align=center, rounded corners=2pt},
    boxfinal/.style={draw, fill=spaet, text=white, minimum width=2.5cm, minimum height=1.2cm, align=center, rounded corners=2pt},
    arrow/.style={-{Stealth[length=3mm]}, thick, color=spaet}
]
\node[box] (start) {A1\\{\scriptsize Beginn}};
\node[box, right=of start] (a2) {A2\\{\scriptsize Ende Klasse 10}};
\node[box, right=of a2] (a2p) {A2+\\{\scriptsize Ende Klasse 11}};
\node[boxfinal, right=of a2p] (b1) {B1(+)\\{\scriptsize Abitur}};

\draw[arrow] (start) -- (a2);
\draw[arrow] (a2) -- (a2p);
\draw[arrow] (a2p) -- (b1);

\node[below=0.5cm of start, align=center] {\scriptsize 160 Stunden\\{\tiny (4 Wochenstunden $\times$ 40 Wochen)}};
\node[below=0.5cm of a2, align=center] {\scriptsize 160 Stunden};
\node[below=0.5cm of a2p, align=center] {\scriptsize 160 Stunden};
\node[below=0.5cm of b1, align=center] {\scriptsize $\Sigma$ 480 Stunden};
\end{tikzpicture}
\end{center}

\hinweis{
\textbf{Quelle:} Handreichung Sprechprüfung (2020), Seite 3:\\
\textit{``In der Einführungsphase neu begonnene Fremdsprachen orientieren sich am Ende der Einführungsphase am Niveau A2 und am Ende der Qualifikationsphase am Zielniveau B1(+).''}
}

\subsubsection{Vergleich mit fortgeführter Fremdsprache}

\begin{center}
\small
\begin{tabular}{|l|c|c|c|c|}
\hline
\rowcolor{lightgray}
\textbf{Phase} & \multicolumn{2}{c|}{\textbf{Fortgeführte Fremdsprache}} & \multicolumn{2}{c|}{\textbf{Neu begonnene Fremdsprache}} \\
\cline{2-5}
\rowcolor{lightgray}
& Niveau & Wochenstunden & Niveau & Wochenstunden \\
\hline
Einführungsphase (Klasse 10) & B1 & 3 oder 5 & \textcolor{spaet}{\textbf{A2}} & \textcolor{spaet}{\textbf{4}} \\
\hline
Qualifikationsphase 1 (Klasse 11) & B1+ & 3 oder 5 & \textcolor{spaet}{\textbf{A2+}} & \textcolor{spaet}{\textbf{4}} \\
\hline
Qualifikationsphase 2 (Klasse 12) & B2 & 3 oder 5 & \textcolor{spaet}{\textbf{B1(+)}} & \textcolor{spaet}{\textbf{4}} \\
\hline
\textbf{Abitur} & \textbf{B2} & & \textcolor{spaet}{\textbf{B1(+)}} & $\Sigma\approx$\textbf{480 Stunden} \\
\hline
\end{tabular}
\end{center}

\subsection{Die vier verbindlichen Themenfelder}

Der Rahmenplan Spanisch Qualifikationsphase M-V (2019) definiert \textbf{vier verbindliche Themenfelder}, die für alle Kurse gelten -- auch für die neu begonnene Fremdsprache. Die Tiefe und Komplexität wird jedoch dem Sprachniveau angepasst.

\begin{center}
\small
\begin{tabular}{|c|p{6cm}|c|c|}
\hline
\rowcolor{lightgray}
\textbf{Nr.} & \textbf{Themenfeld} & \textbf{Fortgeführt} & \textbf{Spätbeginner} \\
\hline
1 & El individuo y la sociedad (Individuum und Gesellschaft) & 45 Stunden & \textcolor{spaet}{30--35 Stunden} \\
\hline
2 & La identidad nacional y la diversidad cultural (Nationale Identität und kulturelle Vielfalt) & 45 Stunden & \textcolor{spaet}{30--35 Stunden} \\
\hline
3 & Aspectos actuales de la política y la sociedad (Aktuelle Aspekte aus Politik und Gesellschaft) & 30 Stunden & \textcolor{spaet}{20--25 Stunden} \\
\hline
4 & Desafíos modernos (Moderne Herausforderungen) & 30 Stunden & \textcolor{spaet}{20--25 Stunden} \\
\hline
\rowcolor{lightspaet}
& \textbf{GESAMT} & \textbf{150 Stunden} & \textcolor{spaet}{\textbf{100--120 Stunden}} \\
\hline
\end{tabular}
\end{center}

\quelle{Rahmenplan Spanisch Qualifikationsphase M-V (2019), Seiten 16--19}

\wichtig{
\textbf{Verbindlich vs. empfohlen:}\\[4pt]
Der Rahmenplan unterscheidet zwischen \textbf{linker Spalte} (verbindliche Inhalte) und \textbf{rechter Spalte} (Hinweise und Anregungen). Die linke Spalte \textbf{muss} behandelt werden, die rechte Spalte bietet Beispiele zur Schwerpunktsetzung.\\[4pt]
\textit{``Die [...] benannten Themen und Ziele füllen ca. \textbf{80\%} der zur Verfügung stehenden Unterrichtszeit. [...] 20\% stehen für eigene Unterrichtsgestaltung zur Verfügung.''} (Rahmenplan, Seite 1)
}

\newpage

%% ============================================================================
%% TEIL 2: REGULÄRER UNTERRICHT NACH KLASSENSTUFE
%% ============================================================================

\section*{TEIL 2: REGULÄRER UNTERRICHT NACH KLASSENSTUFE}
\addcontentsline{toc}{section}{TEIL 2: REGULÄRER UNTERRICHT NACH KLASSENSTUFE}

\subsection{Klasse 10 (Einführungsphase)}

\subsubsection{Überblick}

\begin{center}
\small
\begin{tabular}{|p{4cm}|p{9cm}|}
\hline
\rowcolor{lightspaet}
\textbf{Zielniveau} & A2 (Gemeinsamer Europäischer Referenzrahmen) \\
\hline
\textbf{Wochenstunden} & 4 \\
\hline
\textbf{Jahresstunden} & ca. 160 (4 $\times$ 40 Wochen) \\
\hline
\textbf{Fokus} & Grundlagen aufbauen, kommunikative Basiskompetenz entwickeln \\
\hline
\textbf{Klausuren} & 2 pro Jahr, mindestens 45 Minuten (90 Minuten bei Aufsätzen) \\
\hline
\textbf{Sprechprüfung} & Optional (empfohlen: noch nicht in Klasse 10) \\
\hline
\textbf{Lehrwerk} & A\_tope.com: Unidades 1--4 \\
\hline
\end{tabular}
\end{center}

\subsubsection{Thematische Inhalte (Halbjahresplanung)}

\textbf{Halbjahr 10.1 (August--Januar):}
\begin{itemize}[noitemsep]
\item Wochen 1--4: Grundlagen (Aussprache, Alphabet, Basisgrammatik)
\item Wochen 5--8: Yo y mi entorno (sich vorstellen, Familie)
\item Wochen 9--12: La vida cotidiana (Alltag, Uhrzeit, Aktivitäten)
\item Wochen 13--16: Mi colegio y mis amigos + \textbf{Klausur 1}
\end{itemize}

\textbf{Halbjahr 10.2 (Februar--Juli):}
\begin{itemize}[noitemsep]
\item Wochen 1--4: España: geografía básica
\item Wochen 5--8: Tradiciones y fiestas
\item Wochen 9--12: De compras y en el restaurante
\item Wochen 13--16: Wiederholung + \textbf{Klausur 2}
\end{itemize}

\subsubsection{Grammatik (Jahr 1)}

\begin{center}
\small
\begin{tabular}{|l|p{8cm}|}
\hline
\rowcolor{lightgray}
\textbf{Struktur} & \textbf{Beschreibung} \\
\hline
Presente de indicativo & Regelmäßige Verben + wichtige unregelmäßige (ser, estar, tener, ir, hacer) \\
\hline
Artículos, género, número & Bestimmte/unbestimmte Artikel, Genus und Numerus \\
\hline
Possessivbegleiter & mi, tu, su, nuestro, vuestro, su \\
\hline
Pretérito perfecto & Einführung (Perfekt mit haber + Partizip) \\
\hline
Futuro próximo & ir a + Infinitiv \\
\hline
Präpositionen & Grundlegende Ortsangaben (en, a, de, con, para, por) \\
\hline
\end{tabular}
\end{center}

\subsubsection{Klausuren in Klasse 10}

\begin{center}
\fcolorbox{darkgray}{lightgray}{\parbox{0.9\textwidth}{
\small
\textbf{Anzahl:} 2 Klausuren pro Schuljahr (1 pro Halbjahr)\\
\textbf{Dauer:} Mindestens 45 Minuten; bei Aufsätzen mindestens 90 Minuten\\
\textbf{Niveau:} A2\\
\textbf{Bewertungsbogen:} A2-Bogen (falls Sprechprüfung)\\
\textbf{Ersetzung:} 1 Klausur pro Jahr kann durch eine Sprechprüfung ersetzt werden\\[4pt]
\quelle{Oberstufen- und Abiturprüfungsverordnung M-V, Handreichung Sprechprüfung (2020)}
}}
\end{center}

\subsubsection{Lehrwerk-Nutzung: A\_tope.com in Klasse 10}

\begin{center}
\small
\begin{tabular}{|l|c|p{6cm}|}
\hline
\rowcolor{lightgray}
\textbf{Einheit} & \textbf{Seiten} & \textbf{Thema} \\
\hline
Unidad 1: Hola, ¿qué tal? & 8--23 & Begrüßung, Vorstellen, Zahlen \\
\hline
Unidad 2: Gente joven & 24--41 & Familie, Alter, Beschreibungen \\
\hline
Unidad 3: Mi mundo & 42--59 & Alltag, Schule, Freizeit \\
\hline
Unidad 4: ¡Buen viaje! & 60--77 & Reisen, Pläne, Zukunft \\
\hline
\end{tabular}
\end{center}

\hinweis{
\textbf{Deckung Rahmenplan:} A\_tope.com deckt die Grundlagen für Klasse 10 zu \textbf{ca. 90\%} ab. Das Lehrwerk ist für diese Klassenstufe \textbf{ausreichend} ohne Ergänzungen.
}

\newpage

\subsection{Klasse 11 (Qualifikationsphase 1)}

\subsubsection{Überblick}

\begin{center}
\small
\begin{tabular}{|p{4cm}|p{9cm}|}
\hline
\rowcolor{lightspaet}
\textbf{Zielniveau} & A2+ $\rightarrow$ B1 (Progression) \\
\hline
\textbf{Wochenstunden} & 4 \\
\hline
\textbf{Jahresstunden} & ca. 160 \\
\hline
\textbf{Themenfelder} & 1 (El individuo) + 2 (Identidad nacional) -- vereinfacht \\
\hline
\textbf{Klausuren} & 2 pro Jahr, mindestens 90 Minuten \\
\hline
\textbf{Sprechprüfung} & \textbf{Empfohlen} (kann 1 Klausur ersetzen) \\
\hline
\textbf{Lehrwerk} & A\_tope.com: Unidades 5--8 + Module \\
\hline
\end{tabular}
\end{center}

\subsubsection{Thematische Inhalte (Halbjahresplanung)}

\textbf{Halbjahr 11.1: El individuo y la sociedad (Themenfeld 1)}
\begin{itemize}[noitemsep]
\item Wochen 1--4: Jóvenes y su mundo (Jugendliche und ihre Welt)
\item Wochen 5--8: Relaciones y emociones (Beziehungen und Emotionen)
\item Wochen 9--12: Perspectivas futuras (Beruf, Studium)
\item Wochen 13--16: \textbf{Klausur 1} + Wiederholung
\end{itemize}

\textbf{Halbjahr 11.2: La identidad nacional (Themenfeld 2)}
\begin{itemize}[noitemsep]
\item Wochen 1--4: España: regiones y cultura (Spanien: Regionen und Kultur)
\item Wochen 5--8: Latinoamérica: introducción (Lateinamerika: Einführung)
\item Wochen 9--12: Diversidad cultural (Kulturelle Vielfalt)
\item Wochen 13--16: \textbf{Sprechprüfung} (B1) oder \textbf{Klausur 2}
\end{itemize}

\subsubsection{Grammatik (Jahr 2)}

\begin{center}
\small
\begin{tabular}{|l|p{8cm}|}
\hline
\rowcolor{lightgray}
\textbf{Struktur} & \textbf{Beschreibung} \\
\hline
Pretérito indefinido & Vollständige Einführung (regelmäßig + unregelmäßig) \\
\hline
Pretérito imperfecto & Formen und Verwendung \\
\hline
Kontrast indefinido/imperfecto & Unterscheidung in Erzählungen \\
\hline
Subjuntivo presente & Einführung (que + Subjuntivo nach Wunsch, Aufforderung) \\
\hline
Condicional simple & Konditional für höfliche Bitten, Wünsche \\
\hline
Relativsätze & que, donde, quien \\
\hline
\end{tabular}
\end{center}

\subsubsection{Klausuren in Klasse 11}

\begin{center}
\fcolorbox{darkgray}{lightgray}{\parbox{0.9\textwidth}{
\small
\textbf{Anzahl:} 2 Klausuren pro Schuljahr (1 pro Halbjahr)\\
\textbf{Dauer:} Mindestens 90 Minuten\\
\textbf{Niveau:} A2+/B1\\
\textbf{Bewertungsbogen:} A2+ oder B1 (je nach Halbjahr)\\
\textbf{Ersetzung:} 1 Klausur pro Jahr kann durch eine Sprechprüfung ersetzt werden\\[4pt]
\textbf{Empfehlung:} Sprechprüfung in Klasse 11 durchführen (B1-Bogen), um Entlastung in Klasse 12 zu schaffen.\\[4pt]
\quelle{Oberstufen- und Abiturprüfungsverordnung M-V, Handreichung Sprechprüfung (2020)}
}}
\end{center}

\subsubsection{Lehrwerk-Nutzung: A\_tope.com in Klasse 11}

\begin{center}
\small
\begin{tabular}{|l|c|p{6cm}|}
\hline
\rowcolor{lightgray}
\textbf{Einheit} & \textbf{Seiten} & \textbf{Thema} \\
\hline
Unidad 5: ¡Bienvenidos a Madrid! & 78--95 & Madrid, Stadtleben \\
\hline
Unidad 6: Peru: ayer y hoy & 96--113 & Peru, Geschichte, Kultur \\
\hline
Unidad 7: Buscando trabajo & 114--131 & Berufswelt, Bewerbung \\
\hline
Módulo: Vivir la diversidad & 126--127 & Toleranz, Vielfalt in Spanien \\
\hline
\end{tabular}
\end{center}

\hinweis{
\textbf{Deckung Rahmenplan:} A\_tope.com deckt die Themenfelder 1 und 2 zu \textbf{ca. 75\%} ab.\\
\textbf{Lücke:} Lateinamerika wird nur durch Peru repräsentiert. Für eine breitere Perspektive sind Ergänzungen sinnvoll (z.B. Mexiko, Argentinien).
}

\newpage

\subsection{Klasse 12 (Qualifikationsphase 2 und Abitur-Vorbereitung)}

\subsubsection{Überblick}

\begin{center}
\small
\begin{tabular}{|p{4cm}|p{9cm}|}
\hline
\rowcolor{lightspaet}
\textbf{Zielniveau} & B1(+) (Abiturniveau) \\
\hline
\textbf{Wochenstunden} & 4 \\
\hline
\textbf{Jahresstunden} & ca. 160 \\
\hline
\textbf{Themenfelder} & 3 (Aspectos actuales) + 4 (Desafíos modernos) \\
\hline
\textbf{Klausuren} & 2 pro Jahr, mindestens 90 Minuten (letzte: abiturnah) \\
\hline
\textbf{Sprechprüfung} & \textbf{Pflicht}, falls noch nicht erfolgt \\
\hline
\textbf{Lehrwerk} & A\_tope.com: Unidad 8 + Module -- \textbf{Ergänzungen erforderlich!} \\
\hline
\end{tabular}
\end{center}

\subsubsection{Thematische Inhalte (Halbjahresplanung)}

\textbf{Halbjahr 12.1: Aspectos actuales de la política y la sociedad (Themenfeld 3)}
\begin{itemize}[noitemsep]
\item Wochen 1--4: \textbf{Migración y sus causas} (Migration und ihre Ursachen) -- \textcolor{spanisch}{\textbf{Ergänzung erforderlich!}}
\item Wochen 5--8: Globalización (Globalisierung)
\item Wochen 9--12: \textbf{Klausur 1} (abiturnah)
\item Wochen 13--16: Wiederholung Themenfeld 1+2
\end{itemize}

\textbf{Halbjahr 12.2: Desafíos modernos + Abitur (Themenfeld 4)}
\begin{itemize}[noitemsep]
\item Wochen 1--4: \textbf{Medio ambiente} (Umwelt) -- \textcolor{spanisch}{\textbf{Ergänzung erforderlich!}}
\item Wochen 5--6: \textbf{Tecnología y redes sociales} (Technologie und soziale Medien)
\item Wochen 7--10: \textbf{Abitur-Vorbereitung} (alle Themenfelder) $\rightarrow$ siehe Anhang A
\item Wochen 11+: Abiturprüfung
\end{itemize}

\subsubsection{Grammatik (Jahr 3)}

\begin{center}
\small
\begin{tabular}{|l|p{6cm}|c|}
\hline
\rowcolor{lightgray}
\textbf{Struktur} & \textbf{Beschreibung} & \textbf{Niveau} \\
\hline
Subjuntivo presente & Vertiefung (nach Emotionen, Meinungen) & B1 \\
\hline
Pluscuamperfecto & Vorvergangenheit & B2 (optional) \\
\hline
Voz pasiva & Passivkonstruktionen & B2 (optional) \\
\hline
Indirekte Rede & Wiedergabe von Gesagtem & B1 \\
\hline
Konnektoren & Textverknüpfung (sin embargo, además, por lo tanto) & B1 \\
\hline
\end{tabular}
\end{center}

\hinweis{
\textbf{Hinweis für Spätbeginner:} Pluscuamperfecto und Passiv sind B2-Strukturen und für das Zielniveau B1(+) \textbf{nicht zwingend erforderlich}. Der Fokus sollte auf der sicheren Beherrschung der B1-Grammatik liegen.
}

\subsubsection{Klausuren in Klasse 12}

\begin{center}
\fcolorbox{darkgray}{lightgray}{\parbox{0.9\textwidth}{
\small
\textbf{Anzahl:} 2 Klausuren pro Schuljahr (nur 1 Klausur im Abitur-Semester)\\
\textbf{Dauer:} Mindestens 90 Minuten\\
\textbf{Niveau:} B1\\
\textbf{Besonderheit:} Die letzte Klausur vor dem Abitur sollte \textbf{abiturnah} gestaltet sein (Format, Dauer, Anforderungsbereiche)\\[4pt]
\textbf{Sprechprüfung:} Falls noch nicht in Klasse 10 oder 11 erfolgt, \textbf{muss} mindestens eine Sprechprüfung in der gesamten Oberstufe durchgeführt werden.\\[4pt]
\quelle{Oberstufen- und Abiturprüfungsverordnung M-V}
}}
\end{center}

\subsubsection{Lehrwerk-Nutzung: A\_tope.com in Klasse 12}

\wichtig{
\textbf{KRITISCH: Ergänzungen zwingend erforderlich!}\\[4pt]
Das Lehrwerk A\_tope.com deckt die Themenfelder 3 und 4 \textbf{unzureichend} ab:
\begin{itemize}[noitemsep]
\item \textbf{Migration:} 0\% Deckung (kein inhaltlicher Treffer)
\item \textbf{Umwelt:} ca. 15\% (nur ``turismo sostenible'' in Unidad 8)
\item \textbf{Digitalisierung:} nur Vokabular, keine kritische Reflexion
\end{itemize}
}

\begin{center}
\small
\begin{tabular}{|l|c|p{5cm}|}
\hline
\rowcolor{lightgray}
\textbf{Einheit} & \textbf{Seiten} & \textbf{Thema} \\
\hline
Unidad 8: Bienvenidos a Andalucía & 132--149 & Andalusien, Tourismus, Flamenco \\
\hline
Módulo: El comercio justo & 128--129 & Fairer Handel \\
\hline
Módulo: Una decisión importante & 130--131 & Emigration (Spanier nach Deutschland) \\
\hline
Módulo: ¿Te comunicas? & 132--133 & Handy-Vokabular \\
\hline
\end{tabular}
\end{center}

\subsubsection{Erforderliche Ergänzungen für Klasse 12}

\begin{center}
\small
\begin{tabular}{|p{4cm}|p{5cm}|p{4cm}|}
\hline
\rowcolor{lightgray}
\textbf{Verbindlicher Inhalt} & \textbf{Ergänzung} & \textbf{Quelle} \\
\hline
\rowcolor{lightred}
\textbf{Movimientos migratorios} (Migration) & Lektüre \textit{Abdel} (Enrique Páez) + Handreichung Themenfeld 3 & Literaturempfehlungen M-V \\
\hline
\rowcolor{lightred}
\textbf{Tendencias ecológicas} (Umwelt) & Handreichung Themenfeld 4 + Aktualtexte (cambio climático, contaminación) & Rahmenplan Seite 19 \\
\hline
\rowcolor{lightred}
\textbf{Desafíos de la era digital} (Digitalisierung) & Aktualtexte zur kritischen Medienreflexion (Chancen und Risiken) & Rahmenplan Seite 19 \\
\hline
Globalización & A\_tope Módulo ``El comercio justo'' + Aktualtexte & -- \\
\hline
\end{tabular}
\end{center}

\hinweis{
\textbf{Literatur \textit{Abdel}:} 16-jähriger Protagonist aus Marokko, Thema illegale Immigration. Niveau A2/B1 (sprachlich zugänglich). Chronologische Erzählung. \textbf{Praktisch unverzichtbar} für Themenfeld 3, da A\_tope Migration nicht behandelt.\\[4pt]
\textbf{Ist Literatur verbindlich?} Nein -- die Literaturempfehlungen sind ein \textit{Begleitdokument}, keine Vorgabe. Aber: Da das Thema Migration im Lehrwerk fehlt, ist \textit{Abdel} die praktisch einzige zugängliche Option für Spätbeginner.
}

\newpage

%% ============================================================================
%% TEIL 3: PRÜFUNGSFORMATE
%% ============================================================================

\section*{TEIL 3: PRÜFUNGSFORMATE}
\addcontentsline{toc}{section}{TEIL 3: PRÜFUNGSFORMATE}

\subsection{Klausuren in der Oberstufe}

\subsubsection{Anzahl und Dauer}

\begin{center}
\small
\begin{tabular}{|l|c|c|c|}
\hline
\rowcolor{lightgray}
\textbf{Phase} & \textbf{Klausuren/Jahr} & \textbf{Mindestdauer} & \textbf{Niveau} \\
\hline
Einführungsphase (Klasse 10) & 2 & 45 Minuten (90 bei Aufsätzen) & A2 \\
\hline
Qualifikationsphase (Klasse 11) & 2 & 90 Minuten & A2+/B1 \\
\hline
Qualifikationsphase (Klasse 12) & 2 (nur 1 im Abitur-Semester) & 90 Minuten & B1 \\
\hline
\end{tabular}
\end{center}

\quelle{Oberstufen- und Abiturprüfungsverordnung M-V}

\subsubsection{Ersetzung durch Sprechprüfung}

\begin{center}
\begin{tikzpicture}[
    node distance=1.5cm,
    box/.style={draw, fill=lightspaet, minimum width=4cm, minimum height=1cm, align=center, rounded corners=2pt},
    decision/.style={draw, diamond, fill=lightgray, minimum width=2cm, aspect=2, align=center},
    arrow/.style={-{Stealth[length=3mm]}, thick}
]
\node[box] (start) {Pro Schuljahr:\\max. 1 Klausur ersetzbar};
\node[decision, below=of start] (check) {Sprechprüfung\\durchgeführt?};
\node[box, below left=1.5cm and 0.5cm of check] (yes) {Ersetzt 1 Klausur\\im Zeugnis};
\node[box, below right=1.5cm and 0.5cm of check] (no) {Reguläre Klausur\\schreiben};

\draw[arrow] (start) -- (check);
\draw[arrow] (check) -- node[left] {Ja} (yes);
\draw[arrow] (check) -- node[right] {Nein} (no);
\end{tikzpicture}
\end{center}

\wichtig{
\textbf{Mindestens 1 Sprechprüfung verpflichtend!}\\[4pt]
Laut § 22 Absatz 5 der Oberstufen- und Abiturprüfungsverordnung M-V muss \textbf{jede Schülerin und jeder Schüler} mindestens einmal in der gesamten Oberstufe (Klassen 10--12) eine komplexe Leistungsermittlung im Kompetenzbereich Sprechen absolvieren.
}

\subsection{Sprechprüfung (verbindlich)}

\subsubsection{Format: Drei Teile}

\begin{center}
\small
\begin{tabular}{|c|p{3.5cm}|c|p{6cm}|}
\hline
\rowcolor{lightgray}
\textbf{Teil} & \textbf{Name} & \textbf{Dauer} & \textbf{Beschreibung} \\
\hline
1 & Interview & 2 Minuten & Aufwärmen, persönliche Fragen -- \textbf{keine Bewertung!} \\
\hline
2 & Monolog & 4 Minuten/Schüler & Impulsgesteuerte Präsentation (Bild, Text, Grafik) \\
\hline
3 & Dialog & 5 Minuten/Schüler & Diskussion oder Rollenspiel zu einem Thema \\
\hline
\multicolumn{2}{|c|}{\textbf{Gesamt}} & \textbf{20 Minuten} & Tandem-Prüfung (2 Schüler gleichzeitig) \\
\hline
\end{tabular}
\end{center}

\subsubsection{Bewertung}

\begin{center}
\small
\begin{tabular}{|p{6cm}|c|}
\hline
\rowcolor{lightgray}
\textbf{Kriterium} & \textbf{Gewichtung} \\
\hline
\textbf{Sprachliche Leistung} & \textbf{60\%} \\
\quad Aussprache und Intonation & \\
\quad Wortschatz und Ausdruck & \\
\quad Grammatische Korrektheit & \\
\quad Flüssigkeit und Interaktion & \\
\hline
\textbf{Inhaltliche Leistung} & \textbf{40\%} \\
\quad Aufgabenerfüllung & \\
\quad Argumentation & \\
\quad Themenentwicklung & \\
\hline
\end{tabular}
\end{center}

\subsubsection{Bewertungsbögen nach Niveau}

\begin{center}
\small
\begin{tabular}{|l|c|c|}
\hline
\rowcolor{lightgray}
\textbf{Phase} & \textbf{Bewertungsbogen} & \textbf{Quelle} \\
\hline
Einführungsphase (Klasse 10) -- Spätbeginner & \textcolor{spaet}{\textbf{A2}} & Handreichung Sprechprüfung, Anlage II \\
\hline
Qualifikationsphase (Klasse 11/12) -- Spätbeginner & \textcolor{spaet}{\textbf{B1}} & Handreichung Sprechprüfung, Anlage II \\
\hline
Fortgeführte Fremdsprache & B2 & Handreichung Sprechprüfung, Anlage II \\
\hline
\end{tabular}
\end{center}

\subsubsection{Kombination Sprechen + Hörverstehen}

Falls die Sprechprüfung mit einer Hör-/Hörsehverstehensaufgabe kombiniert wird (§ 22 Absatz 5 der Oberstufen- und Abiturprüfungsverordnung):

\begin{center}
\small
\begin{tabular}{|l|c|}
\hline
\rowcolor{lightgray}
\textbf{Kompetenzbereich} & \textbf{Gewichtung} \\
\hline
Sprechen & \textbf{80\%} \\
\hline
Hör-/Hörsehverstehen & 20\% \\
\hline
\end{tabular}
\end{center}

\subsubsection{Timing-Empfehlung}

\hinweis{
\textbf{Empfehlung der Handreichung Sprechprüfung (2020), Seite 6:}\\[4pt]
\textit{``Es wird empfohlen, im \textbf{zweiten Schulhalbjahr} die zweiten bzw. dritten Fremdsprachen zu überprüfen.''}\\[4pt]
\textbf{Grund:} Im 4. Halbjahr korrigieren Englisch-Lehrkräfte die schriftlichen Abiturprüfungen. Durch die Verlagerung der Spanisch-Sprechprüfung ins 2. Halbjahr wird eine Korrekturentlastung erreicht.\\[4pt]
\textbf{Praktische Empfehlung:} Sprechprüfung in \textbf{Klasse 11, 2. Halbjahr} durchführen (B1-Bogen). So bleibt in Klasse 12 mehr Zeit für die Abitur-Vorbereitung.
}

\newpage

%% ============================================================================
%% ANHANG A: ABITURPRÜFUNG
%% ============================================================================

\section*{ANHANG A: ABITURPRÜFUNG}
\addcontentsline{toc}{section}{ANHANG A: ABITURPRÜFUNG}

\subsection{A1. Schriftliches Abitur}

\subsubsection{Rahmendaten}

\begin{center}
\small
\begin{tabular}{|p{5cm}|p{8cm}|}
\hline
\rowcolor{lightspaet}
\textbf{Dauer} & 180 Minuten (nur grundlegendes Niveau, kein Leistungskurs!) \\
\hline
\textbf{Niveau} & B1(+) \\
\hline
\textbf{Prüfungsfach} & Wählbar als P4 (Kombination schriftlich+mündlich) \\
\hline
\textbf{Aufgabenformate} & Textaufgabe (Pflicht), Leseverstehen, Hörverstehen (optional), Sprachmittlung (optional) \\
\hline
\end{tabular}
\end{center}

\subsubsection{Anforderungsbereiche (Spätbeginner)}

\begin{center}
\small
\begin{tabular}{|l|c|c|p{5cm}|}
\hline
\rowcolor{lightgray}
\textbf{Anforderungsbereich} & \textbf{Spätbeginner} & \textbf{Fortgeführt} & \textbf{Typische Operatoren} \\
\hline
I (Reproduktion) & \textcolor{spaet}{\textbf{25\%}} & 20\% & nenne, beschreibe, gib an \\
\hline
II (Reorganisation) & \textcolor{spaet}{\textbf{50\%}} & 45\% & erkläre, vergleiche, analysiere \\
\hline
III (Transfer) & \textcolor{spaet}{\textbf{25\%}} & 35\% & beurteile, diskutiere, entwickle \\
\hline
\end{tabular}
\end{center}

\hinweis{
\textbf{Für Spätbeginner gilt:} Mehr Gewicht auf Anforderungsbereich I (25\% statt 20\%) und weniger auf Anforderungsbereich III (25\% statt 35\%). Dies berücksichtigt das niedrigere Sprachniveau B1(+) gegenüber B2.
}

\subsection{A2. Mündliche Abiturprüfung (Paarprüfung)}

\subsubsection{Rechtliche Grundlage}

\begin{itemize}[noitemsep]
\item § 25 Absatz 8 der Oberstufen- und Abiturprüfungsverordnung M-V (Kombination P4)
\item § 38 Absatz 6 der Oberstufen- und Abiturprüfungsverordnung M-V (mündliche Prüfung P5)
\item Verwaltungsvorschrift Paarprüfung vom 03.12.2025
\item Handreichung Paarprüfung Abitur (2025)
\end{itemize}

\subsubsection{Anwendungsfälle}

\begin{center}
\small
\begin{tabular}{|l|p{10cm}|}
\hline
\rowcolor{lightgray}
\textbf{Prüfungsfach} & \textbf{Beschreibung} \\
\hline
\textbf{P5} (mündliche Abiturprüfung) & Vollständig mündlich, Paarprüfung möglich \\
\hline
\textbf{P4} (Kombination) & Verkürzte schriftliche Prüfung + mündliche Prüfung \\
\hline
\end{tabular}
\end{center}

\subsubsection{Format der Paarprüfung}

\begin{center}
\small
\begin{tabular}{|c|p{3.5cm}|c|p{6cm}|}
\hline
\rowcolor{lightgray}
\textbf{Teil} & \textbf{Name} & \textbf{Dauer} & \textbf{Beschreibung} \\
\hline
1 & Interview & 2 Minuten & Kurzer Einstieg, Auflockerung -- \textbf{keine Bewertung!} \\
\hline
2 & Monolog & 4 Minuten/Prüfling & Längerer Gesprächsbeitrag auf Basis eines Impulses \\
\hline
3 & Dialog & 5 Minuten/Prüfling & Interaktion zwischen beiden Prüflingen \\
\hline
\multicolumn{2}{|c|}{\textbf{Gesamt}} & \textbf{ca. 30 Minuten} & Vorbereitungszeit: 20 Minuten (beide Prüflinge gemeinsam) \\
\hline
\end{tabular}
\end{center}

\subsubsection{Bewertung}

\begin{itemize}[noitemsep]
\item \textbf{Individuelle Bewertung} jedes Prüflings (keine Beeinflussung durch Partner)
\item 60\% sprachlich / 40\% inhaltlich
\item \textbf{Bewertungsbogen für Spätbeginner:} B1 (Handreichung Paarprüfung, Anlage I)
\item Impulsmaterial von Prüfern gewählt; beide Prüflinge erhalten unterschiedliche Teilimpulse
\item Im Dialog wird eigenständige Interaktion erwartet; Prüfer greifen nur bei Bedarf ein
\end{itemize}

\subsection{A3. Zusammenfassung: Prüfungen für Spätbeginner}

\begin{center}
\small
\begin{tabular}{|p{3.5cm}|c|c|c|p{4cm}|}
\hline
\rowcolor{lightspaet}
\textbf{Prüfung} & \textbf{Anzahl} & \textbf{Dauer} & \textbf{Niveau} & \textbf{Besonderheit} \\
\hline
Klausur Klasse 10 & 2/Jahr & 45--90 Min & A2 & 1$\times$/Jahr ersetzbar \\
\hline
Klausur Klasse 11--12 & 2/Jahr & 90 Min & A2+/B1 & 1$\times$/Jahr ersetzbar \\
\hline
\textbf{Sprechprüfung} & \textbf{$\geq$1 ges.} & 20 Min & A2/B1 & \textbf{Pflicht!} \\
\hline
Abitur schriftlich & 1 & 180 Min & B1(+) & Nur grundlegendes Niveau \\
\hline
Abitur mündlich P5 & 1 & 30 Min & B1 & Paarprüfung möglich \\
\hline
Abitur Kombi P4 & 1 & variabel & B1 & Verkürzt schriftl. + mündl. \\
\hline
\end{tabular}
\end{center}

\newpage

%% ============================================================================
%% ANHANG B: LEHRWERK A_TOPE.COM
%% ============================================================================

\section*{ANHANG B: LEHRWERK A\_TOPE.COM}
\addcontentsline{toc}{section}{ANHANG B: LEHRWERK A\_TOPE.COM}

\subsection{B1. Übersicht und Metadaten}

\begin{center}
\small
\begin{tabular}{|p{4cm}|p{9cm}|}
\hline
\rowcolor{lightgray}
\textbf{Merkmal} & \textbf{Details} \\
\hline
\textbf{Titel} & A\_tope.com -- Spanisch Spätbeginner \\
\hline
\textbf{Verlag} & Cornelsen Verlag \\
\hline
\textbf{Erscheinungsjahr} & 2010 (aktualisiert 2017, Druck 2020) \\
\hline
\textbf{Zielgruppe} & Gymnasium, Gesamtschule, Berufsschule, Erwachsenenbildung (bundesweit) \\
\hline
\textbf{Zulassung} & Alle 16 Bundesländer, einschließlich Mecklenburg-Vorpommern \\
\hline
\textbf{Umfang} & 8 Unidades + 4 Module + Anhänge \\
\hline
\end{tabular}
\end{center}

\wichtig{
\textbf{Kernproblem:} A\_tope.com (2010/2017) wurde \textbf{vor} dem Rahmenplan M-V (2019) konzipiert.\\[4pt]
Die \textbf{formale Zulassung} für M-V bedeutet \textbf{nicht}, dass alle verbindlichen Inhalte des M-V-Rahmenplans abgedeckt werden.
}

\subsection{B2. Seitenübersicht nach Themen}

\subsubsection{Unidades (Haupteinheiten)}

\begin{center}
\small
\begin{tabular}{|c|l|c|c|p{4cm}|}
\hline
\rowcolor{lightgray}
\textbf{Nr.} & \textbf{Titel} & \textbf{Seiten} & \textbf{Themenfeld} & \textbf{Grammatik} \\
\hline
1 & Hola, ¿qué tal? & 8--23 & -- (Grundlagen) & Presente, Artikel \\
\hline
2 & Gente joven & 24--41 & 1 & Possessivbegleiter \\
\hline
3 & Mi mundo & 42--59 & 1 & Perfecto \\
\hline
4 & ¡Buen viaje! & 60--77 & 1 & Futuro próximo \\
\hline
5 & ¡Bienvenidos a Madrid! & 78--95 & 2 & Indefinido \\
\hline
6 & Peru: ayer y hoy & 96--113 & 2 & Imperfecto \\
\hline
7 & Buscando trabajo & 114--131 & 1, 3 & Condicional \\
\hline
8 & Bienvenidos a Andalucía & 132--149 & 2, 4 & (Wiederholung) \\
\hline
\end{tabular}
\end{center}

\subsubsection{Module (fakultativ)}

\begin{center}
\small
\begin{tabular}{|l|c|p{5cm}|c|}
\hline
\rowcolor{lightgray}
\textbf{Modul} & \textbf{Seiten} & \textbf{Tatsächlicher Inhalt} & \textbf{Themenfeld} \\
\hline
Vivir la diversidad & 126--127 & Toleranz-Kampagne \textbf{innerhalb Spaniens}, Decálogo gegen Diskriminierung & 2 (nicht 3!) \\
\hline
El comercio justo & 128--129 & 10 Gründe für fairen Handel, Kaufverhalten & 3 (teilweise) \\
\hline
Una decisión importante & 130--131 & Spanischer Ingenieur \textbf{emigriert nach Deutschland} & 3 (aber Emigration!) \\
\hline
¿Te comunicas? & 132--133 & Handy-Aktivitäten: Nachrichten, Fotos, Apps & 4 (nur Vokabular) \\
\hline
\end{tabular}
\end{center}

\hinweis{
\textbf{Wichtiger Befund zum Modul ``Una decisión importante'':}\\[4pt]
Dieses Modul behandelt \textbf{Emigration} (ein Spanier geht wegen Jobmangel nach Deutschland), \textbf{nicht} Immigration nach Spanien. Das Modul deckt somit \textbf{nicht} den verbindlichen Inhalt ``Movimientos migratorios'' ab, der im Rahmenplan die \textit{Analyse und Diskussion der Ursachen und Konsequenzen von Migration} fordert.
}

\subsection{B3. Deckungsmatrix: Rahmenplan vs. Lehrwerk}

\begin{center}
\small
\begin{tabular}{|c|p{4.5cm}|c|c|p{4cm}|}
\hline
\rowcolor{lightgray}
\textbf{Nr.} & \textbf{Themenfeld} & \textbf{Rahmenplan} & \textbf{A\_tope} & \textbf{Befund} \\
\hline
1 & El individuo y la sociedad & \checkmark & \textcolor{success}{\textbf{90\%}} & Unidad 1--4, 7: Jugend, Familie, Beruf \\
\hline
2 & Identidad nacional / Diversidad & \checkmark & \textcolor{warning}{\textbf{70\%}} & Unidad 5, 6, 8 -- Lücke: nur Peru für Lateinamerika \\
\hline
\rowcolor{lightred}
\textbf{3} & \textbf{Aspectos actuales} & \checkmark & \textcolor{spanisch}{\textbf{0\%}} & \textbf{Migration fehlt komplett} \\
\hline
\rowcolor{lightred}
\textbf{4} & \textbf{Desafíos modernos} & \checkmark & \textcolor{spanisch}{\textbf{15\%}} & Umwelt/Digital nur oberflächlich \\
\hline
\end{tabular}
\end{center}

\subsubsection{Detailanalyse Themenfeld 3}

\begin{center}
\small
\begin{tabular}{|p{4cm}|p{5cm}|p{4cm}|}
\hline
\rowcolor{lightgray}
\textbf{Verbindlicher Inhalt} & \textbf{Rahmenplan (Seite 18)} & \textbf{A\_tope.com} \\
\hline
Desarrollos económicos & ``wirtschaftliche Entwicklungen [...] exemplarisch reflektieren'' & Teilweise: Unidad 7 (Berufswelt) \\
\hline
\rowcolor{lightred}
\textbf{Movimientos migratorios} & ``Analyse und Diskussion der \textbf{Ursachen und Konsequenzen von Migration}'' & \textcolor{spanisch}{\textbf{Nicht vorhanden}} \\
\hline
\end{tabular}
\end{center}

\subsubsection{Detailanalyse Themenfeld 4}

\begin{center}
\small
\begin{tabular}{|p{4cm}|p{5cm}|p{4cm}|}
\hline
\rowcolor{lightgray}
\textbf{Verbindlicher Inhalt} & \textbf{Rahmenplan (Seite 19)} & \textbf{A\_tope.com} \\
\hline
\rowcolor{lightred}
\textbf{Tendencias ecológicas} & ``aktuelle Dokumente zum Thema \textbf{Umweltschutz} sowie Auswirkungen verschiedener Tourismuskonzepte'' & \textcolor{spanisch}{\textbf{Nur Randerwähnung}} (``turismo sostenible'' in Unidad 8) \\
\hline
\rowcolor{lightred}
\textbf{Desafíos de la era digital} & ``Mediennutzung [...] im \textbf{Spannungsfeld zwischen Chancen und Risiken}'' & \textcolor{spanisch}{\textbf{Nur Vokabular}} (Handy-Aktivitäten) \\
\hline
\end{tabular}
\end{center}

\subsection{B4. Erforderliche Ergänzungen}

\subsubsection{Nach Lernjahr}

\begin{center}
\small
\begin{tabular}{|l|c|c|p{6cm}|}
\hline
\rowcolor{lightgray}
\textbf{Jahr} & \textbf{Niveau} & \textbf{Deckung} & \textbf{Bewertung} \\
\hline
Jahr 1 (Klasse 10) & A2 & \textcolor{success}{\textbf{90\%}} & A\_tope ausreichend \\
\hline
Jahr 2 (Klasse 11) & A2+/B1 & \textcolor{warning}{\textbf{75\%}} & A\_tope + Module ausreichend \\
\hline
\rowcolor{lightred}
\textbf{Jahr 3 (Klasse 12)} & B1(+) & \textcolor{spanisch}{\textbf{40\%}} & \textbf{Ergänzungen erforderlich} \\
\hline
\end{tabular}
\end{center}

\subsubsection{Konkrete Ergänzungen für Jahr 3}

\begin{center}
\small
\begin{tabular}{|p{4cm}|p{5cm}|p{4cm}|}
\hline
\rowcolor{lightgray}
\textbf{Thema} & \textbf{Empfohlene Ergänzung} & \textbf{Quelle} \\
\hline
\textbf{Migration} & Lektüre \textit{Abdel} (Enrique Páez) + Handreichung Themenfeld 3 & Literaturempfehlungen M-V \\
\hline
\textbf{Umwelt} & Handreichung Themenfeld 4 + Aktualtexte (cambio climático, contaminación) & Institut für Qualitätsentwicklung \\
\hline
\textbf{Digitalisierung} & Aktualtexte zur kritischen Medienreflexion (Chancen und Risiken) & -- \\
\hline
\textbf{Globalización} & A\_tope Módulo ``El comercio justo'' + Aktualtexte & -- \\
\hline
\end{tabular}
\end{center}

\subsection{B5. Literaturempfehlungen für Spätbeginner}

\subsubsection{Kriterien für geeignete Literatur}

\begin{itemize}[noitemsep]
\item Sprachniveau A2--B1 (nicht höher)
\item Kürzere Texte oder Auszüge
\item Annotierte Ausgaben bevorzugt
\item Klare Handlung, wenig implizite Ebenen
\end{itemize}

\subsubsection{Empfohlene Lektüren}

\begin{center}
\small
\begin{tabular}{|p{3cm}|p{2.5cm}|c|c|p{4cm}|}
\hline
\rowcolor{lightgray}
\textbf{Titel} & \textbf{Autor} & \textbf{Niveau} & \textbf{Themenfeld} & \textbf{Eignung} \\
\hline
\rowcolor{lightspaet}
\textit{Abdel} & Enrique Páez & A2/B1 & 3 & \textcolor{spaet}{\textbf{Ideal für Spätbeginner}} \\
\hline
\textit{La mirada} & Juan Madrid & A2+ & 1 & Kurz, spannend, überschaubar \\
\hline
\textit{Desconexión total} & Lourdes Miquel & A2/B1 & 1 & Gruppenarbeit möglich \\
\hline
\textit{547 amigos} & Lorenzo Silva & B1 & 4 & Aktuelles Thema, kurz \\
\hline
\textit{El desertador} & J.M. Merino & B1 & 2 & Kurzgeschichte, Geschichte \\
\hline
\textit{La Casa} (Auszüge) & Paco Roca & A2+ & 1 & Novela gráfica, visuell \\
\hline
\end{tabular}
\end{center}

\hinweis{
\textbf{Top-Empfehlung: \textit{Abdel} (Enrique Páez)}\\[4pt]
16-jähriger Protagonist aus Marokko. Thema: Illegale Immigration nach Spanien. Sprachlich zugänglich (A2/B1), chronologische Erzählung. Ermöglicht Identifikation durch gleichaltrigen Protagonisten.\\[4pt]
\textbf{Praktisch unverzichtbar} für Themenfeld 3, da das Lehrwerk A\_tope.com das Thema Migration nicht behandelt.\\[4pt]
\textbf{Ist Literatur verbindlich?} Nein -- die Literaturempfehlungen sind ein Begleitdokument, keine Vorgabe. Aber: Ohne alternative Quelle für das Thema Migration ist \textit{Abdel} die einzige zugängliche Option.
}

\newpage

%% ============================================================================
%% ANHANG C: PRAKTISCHE HINWEISE
%% ============================================================================

\section*{ANHANG C: PRAKTISCHE HINWEISE}
\addcontentsline{toc}{section}{ANHANG C: PRAKTISCHE HINWEISE}

\subsection{C1. Jahrgangsübergreifende Kurse}

\subsubsection{Situation an kleineren Schulen}

An Gesamtschulen werden Spätbeginner-Kurse oft \textbf{jahrgangsübergreifend} unterrichtet:
\begin{itemize}[noitemsep]
\item Klassen 10, 11, 12 \textbf{gemeinsam} in einem Kurs
\item Alle mit \textbf{4 Wochenstunden} (konform mit der Oberstufen- und Abiturprüfungsverordnung)
\item Gleicher Unterricht, aber unterschiedliche Prüfungsanforderungen
\end{itemize}

\subsubsection{Ist das zulässig?}

\wichtig{
\textbf{Ja}, wenn folgende Bedingungen erfüllt sind:
\begin{enumerate}[noitemsep]
\item Alle Schülerinnen und Schüler erhalten \textbf{4 Wochenstunden}
\item \textbf{Klausuren differenziert} nach Jahrgang:
\begin{itemize}[noitemsep]
\item Klasse 10: 45--90 Minuten (A2-Niveau)
\item Klasse 11--12: 90 Minuten (A2+/B1-Niveau)
\end{itemize}
\item Bewertung nach \textbf{individuellem Stand}
\item Mindestens 1 Sprechprüfung pro Schüler
\end{enumerate}
Jahrgangsübergreifender Unterricht ist in M-V zulässig, wenn die Vorgaben der Oberstufen- und Abiturprüfungsverordnung eingehalten werden. Die Schulkonferenz muss dies beschließen.
}

\subsubsection{Prüfungsgestaltung bei gemischten Kursen}

\begin{center}
\small
\begin{tabular}{|p{2.5cm}|p{5cm}|p{5cm}|}
\hline
\rowcolor{lightgray}
\textbf{Jahrgang} & \textbf{Klausur} & \textbf{Sprechprüfung} \\
\hline
Klasse 10 (1. Jahr) & 45--90 Minuten, A2-Bogen, 2/Jahr & Optional (Einführung) \\
\hline
Klasse 11 (2. Jahr) & 90 Minuten, A2+/B1-Bogen, 2/Jahr & \textbf{Empfohlen} (ersetzt 1 Klausur) \\
\hline
Klasse 12 (3. Jahr) & 90 Minuten, B1-Bogen, 2/Jahr & Spätestens jetzt \textbf{Pflicht} \\
\hline
\end{tabular}
\end{center}

\subsubsection{Praktische Tipps}

\begin{itemize}[noitemsep]
\item \textbf{Klausuren:} Gleiche Themen, aber unterschiedliche Länge/Komplexität je nach Jahrgang
\item \textbf{Sprechprüfung:} Tandems \textbf{aus gleichem Jahrgang} bilden, oder jahrgangsgemischt mit angepasstem Bewertungsbogen
\item \textbf{Dokumentation:} Jahrgang und Niveau des Gemeinsamen Europäischen Referenzrahmens pro Schüler festhalten
\item \textbf{Binnendifferenzierung:} Effektiv mehrere Lerngruppen in einem Kurs; differenzierte Materialien erforderlich
\end{itemize}

\subsection{C2. Kurszusammensetzung (konkretes Beispiel)}

\begin{center}
\small
\begin{tabular}{|c|c|p{10cm}|}
\hline
\rowcolor{lightgray}
\textbf{Klasse} & \textbf{Schüler} & \textbf{Besonderheiten} \\
\hline
\textbf{10} & 6 & Regelkurs, keine besonderen Förderbedarfe \\
\hline
\textbf{11} & 2 & \textbf{Schülerin A:} Arbeitet mit Klasse 10 (Auslandsaufenthalt, muss Stoff nachholen) \newline \textbf{Schülerin B:} Sehr begabt, arbeitet praktisch mit Klasse 12 \\
\hline
\textbf{12} & 2 & \textbf{Schülerin A:} Lese-Rechtschreib-Schwäche mit Nachteilsausgleich, schwächeres Sprachniveau, bemüht \newline \textbf{Schüler B:} Mittelmäßiges Niveau, inkonsistente Mitarbeit \\
\hline
\rowcolor{lightspaet}
\textbf{Gesamt} & \textbf{10} & Jahrgangsübergreifend, starke Binnendifferenzierung erforderlich \\
\hline
\end{tabular}
\end{center}

\subsubsection{Konsequenzen für die Unterrichtsplanung}

\begin{itemize}[noitemsep]
\item Effektiv \textbf{3 Lerngruppen} in einem Kurs
\item Klasse 10: Grundlagenvermittlung (A1 $\rightarrow$ A2)
\item Klasse 11/12: Vertiefung (A2 $\rightarrow$ B1)
\item Differenzierte Materialien erforderlich
\item \textbf{Nachteilsausgleich für Lese-Rechtschreib-Schwäche beachten} (verlängerte Bearbeitungszeit, größere Schrift, Vorlesen)
\end{itemize}

\subsection{C3. Was ist verbindlich, was ist erlaubt?}

\begin{center}
\small
\begin{tabular}{|p{6cm}|c|p{5cm}|}
\hline
\rowcolor{lightgray}
\textbf{Aspekt} & \textbf{Status} & \textbf{Quelle} \\
\hline
4 Wochenstunden für neu begonnene Fremdsprache & \textbf{Verbindlich} & Oberstufen- und Abiturprüfungsverordnung M-V \\
\hline
4 Themenfelder behandeln & \textbf{Verbindlich} & Rahmenplan Seite 16--19 \\
\hline
Linke Spalte (verbindliche Inhalte) & \textbf{Verbindlich} & Rahmenplan Seite 1 \\
\hline
Rechte Spalte (Hinweise) & Empfohlen & Rahmenplan Seite 1 \\
\hline
Mindestens 1 Sprechprüfung & \textbf{Verbindlich} & § 22 Absatz 5 Oberstufen- und Abiturprüfungsverordnung \\
\hline
Reihenfolge der Themenfelder & Frei wählbar & Rahmenplan Seite 1 \\
\hline
20\% der Unterrichtszeit & Frei gestaltbar & Rahmenplan Seite 1 \\
\hline
Literatur (z.B. \textit{Abdel}) & Empfohlen, nicht Pflicht & Literaturempfehlungen (Begleitdokument) \\
\hline
Jahrgangsübergreifender Unterricht & \textbf{Erlaubt} (mit Schulkonferenzbeschluss) & Schulgesetz M-V \\
\hline
\end{tabular}
\end{center}

\newpage

%% ============================================================================
%% ANHANG: GLOSSAR UND QUELLEN
%% ============================================================================

\section*{GLOSSAR}
\addcontentsline{toc}{section}{GLOSSAR}

\begin{center}
\small
\begin{tabular}{|l|p{10cm}|}
\hline
\rowcolor{lightgray}
\textbf{Begriff} & \textbf{Erklärung} \\
\hline
Anforderungsbereich (I/II/III) & Taxonomie der Aufgabenschwierigkeit: I = Reproduktion, II = Reorganisation/Transfer, III = Reflexion/Problemlösung \\
\hline
Gemeinsamer Europäischer Referenzrahmen & Europäisches Sprachenniveau-System (A1--C2) \\
\hline
Grundkurs & Fortgeführte Fremdsprache mit 3 Wochenstunden \\
\hline
Handreichung Paarprüfung (2025) & Leitfaden für mündliche Abiturprüfungen als Paarprüfung \\
\hline
Handreichung Sprechprüfung (2020) & Leitfaden für Sprechprüfungen in der Oberstufe \\
\hline
Leistungskurs & Fortgeführte Fremdsprache mit 5 Wochenstunden \\
\hline
Neu begonnene Fremdsprache & Spätbeginner: Fremdsprache ab Klasse 10 mit 4 Wochenstunden \\
\hline
Oberstufen- und Abiturprüfungsverordnung M-V & Rechtliche Grundlage für die gymnasiale Oberstufe und das Abitur in Mecklenburg-Vorpommern \\
\hline
P4 / P5 & Prüfungsfach 4 (schriftlich+mündlich kombiniert) / Prüfungsfach 5 (mündlich) \\
\hline
Qualifikationsphase & Klassen 11--12 der gymnasialen Oberstufe \\
\hline
Rahmenplan & Verbindlicher Lehrplan für ein Fach \\
\hline
Themenfeld & Einer der 4 verbindlichen inhaltlichen Schwerpunkte im Rahmenplan \\
\hline
Verwaltungsvorschrift & Ausführungsbestimmung zu einem Gesetz oder einer Verordnung \\
\hline
\end{tabular}
\end{center}

\vspace{1cm}

\section*{QUELLENVERZEICHNIS}
\addcontentsline{toc}{section}{QUELLENVERZEICHNIS}

\begin{enumerate}[noitemsep]
\item \textbf{Rahmenplan für die Qualifikationsphase der gymnasialen Oberstufe -- Spanisch} (2019). Ministerium für Bildung, Wissenschaft und Kultur M-V.
\item \textbf{Oberstufen- und Abiturprüfungsverordnung M-V} (2019). Land Mecklenburg-Vorpommern.
\item \textbf{Handreichung zur komplexen Leistungsermittlung Sprechen} (2020). Institut für Qualitätsentwicklung M-V.
\item \textbf{Handreichung zur Paarprüfung im Abitur} (2025). Institut für Qualitätsentwicklung M-V.
\item \textbf{Kommentierte Literaturempfehlungen Spanisch}. Institut für Qualitätsentwicklung M-V.
\item \textbf{Verwaltungsvorschrift Sprechprüfung} vom 24.07.2020. Mitteilungsblatt des Bildungsministeriums M-V.
\item \textbf{Verwaltungsvorschrift Paarprüfung} vom 03.12.2025. Mitteilungsblatt des Bildungsministeriums M-V.
\item \textbf{A\_tope.com -- Spanisch Spätbeginner} (2010/2017). Cornelsen Verlag.
\end{enumerate}

\vspace{1cm}

\begin{center}
\fcolorbox{spaet}{lightspaet}{\parbox{0.9\textwidth}{
\centering
\small
\textbf{Kompendium Spanisch Spätbeginner -- Mecklenburg-Vorpommern}\\[4pt]
Version 2.0 | Januar 2026\\[4pt]
Methode: Analyse des Rahmenplans + Direktes Lesen des Lehrwerks A\_tope.com (OCR-Volltext, 10.240 Zeilen)\\[4pt]
\textcolor{darkgray}{Lernpfade | Für Lehrkräfte-Anerkennung und Unterrichtsplanung}
}}
\end{center}

\end{document}
